\section{Conceptual Model}\label{conceptual-model}

\textbf{Ontological postulate:} We postulate that the physical universe is a continuous, tensioned three-dimensional brane embedded in a higher-dimensional space. This brane is ontically real; all observable phenomena—particles, forces, spacetime structure—arise from its internal oscillations and patterns of strain. The fourth embedding coordinate represents amplitude deformations of the brane, not an independent traversable dimension. Time is not a geometric coordinate but an external evolution parameter.

Key features of this ontology:
\begin{itemize}
\item The brane is spatially isotropic in its ground state; Lorentz invariance emerges as an effective symmetry from isotropic wave propagation, not as a fundamental property.
\item Amplitude deformations couple to lateral contraction, producing curvature-like effects that mimic gravity. In the weak-field limit, these deformations induce an effective metric tensor $g_{\mu\nu}$ whose Newtonian gravitational potential is sourced by the coarse-grained elastic energy density via a Poisson-like equation.
\item Particles correspond to self-confined solitonic excitations whose topology and internal phase determine charge, spin, and mass.
\item Quantization and collapse-like phenomena arise as nonlinear threshold effects, not stochastic postulates.
\end{itemize}

This section introduces (i) the microscopic substrate model and (ii) the
analytical \emph{layered} interpretation framework that connects substrate waves
to effective gauge structure.
The microscopic equations of motion are always those of the brane embedding
$\mathbf X(x,t)$; all higher-layer objects (spectral envelopes, Berry
connections, effective gauge potentials, etc.) are defined as \emph{diagnostic}
functionals of $\mathbf X$ and $\partial_t\mathbf X$ and never modify the
substrate dynamics (see also Section~\ref{experimental-setting}).

\paragraph{Layered model (overview).}
We distinguish the following layers, which will be referenced throughout the
paper. Each layer is a diagnostic analysis step applied to the substrate; no
layer modifies the dynamics.

\begin{enumerate}
\item \textbf{Substrate layer:} \\
  \emph{Input:} None (fundamental) \\
  \emph{Degrees of freedom:} Real embedding $\mathbf X(x,t) \in \mathbb R^4$ \\
  \emph{Dynamics:} Lagrangian $\mathcal L[\mathbf X, \partial \mathbf X]$ with
  geometric amplitude--lateral coupling \\
  \emph{Output:} Time-evolved configuration $\mathbf X(x,t)$ and velocity
  $\partial_t \mathbf X(x,t)$

\item \textbf{Space--time spectral layer:} \\
  \emph{Input:} Substrate field $\mathbf X(x,t)$ (real-valued) \\
  \emph{Operator:} Local Fourier transform, band-pass filtering in
  $(\mathbf k,\omega)$ space \\
  \emph{Output:} Complex band-isolated envelope $u^A(x,t) \in \mathbb C^4$,
  normalized to $\tilde u(x,t)$

\item \textbf{Berry layer:} \\
  \emph{Input:} Normalized envelope $\tilde u(x,t)$ \\
  \emph{Operator:} Spacetime connection $a_\mu = i \langle \tilde u, \partial_\mu
  \tilde u \rangle$; curvature $f_{\mu\nu} = \partial_\mu a_\nu - \partial_\nu a_\mu$ \\
  \emph{Output:} $U(1)$ gauge connection (or $U(N)$ Wilczek--Zee at degeneracies)

\item \textbf{Electron micro-structure:} \\
  \emph{Input:} Substrate initial conditions (toroidal soliton with W\&vdM
  double-loop, Compton-frequency internal mode) \\
  \emph{Structure:} Localized topological/geometry-driven defects \\
  \emph{Output:} Stable soliton with intrinsic holonomy defect, $4\pi$ periodicity

\item \textbf{Electromagnetism (effective):} \\
  \emph{Input:} Berry connection $a_\mu$ and curvature $f_{\mu\nu}$ \\
  \emph{Operator:} Scaling $A_\mu^{\text{eff}} = (\hbar/q) a_\mu$,
  $F_{\mu\nu}^{\text{eff}} = (\hbar/q) f_{\mu\nu}$ \\
  \emph{Output:} Effective electromagnetic four-potential and field tensor;
  homogeneous Maxwell equations hold by construction (Bianchi)

\item \textbf{Electron (effective):} \\
  \emph{Input:} Non-Abelian holonomy from defect-induced degeneracies \\
  \emph{Structure:} Dirac/spinor as long-wavelength envelope description \\
  \emph{Output:} Spin-$1/2$ representation, effective four-current
  $j^\mu = \bar\psi \gamma^\mu \psi$
\end{enumerate}

The Michelson–Morley null result is naturally explained by Lorentz contraction of matter itself, rendering motion through the medium undetectable. In this view, the observed constancy of light speed reflects a dynamical property of the medium rather than a postulate about spacetime geometry.

\subsection{Continuous Formulation: Pre-tensioned $d$-Brane Embedded in $\mathbb{R}^4$}\label{continuous-formulation}

We model the substrate as a $d$-dimensional material continuum (typically $d=3$) embedded in a Euclidean
$\mathbb{R}^4$. The model has no explicit particles and no additional fields beyond the embedding map itself.
Pre-tension is represented as a \emph{material} (tangential) stress state maintained by boundary conditions; it is
not associated with any particular embedding component.

\paragraph{Kinematics (embedding and induced metric).}
Let $\Omega\subset\mathbb{R}^d$ be material coordinates $x=(x^1,\dots,x^d)$ and let
\begin{equation}
\mathbf{X}:\Omega\times\mathbb{R}\to\mathbb{R}^4,\qquad (x,t)\mapsto \mathbf{X}(x,t)=(X^1,\dots,X^4)
\label{eq:embedding-map}
\end{equation}
be the embedding map. Time $t$ is an external evolution parameter.
The deformation gradient is $F_{Ai}=\partial_i X^A$ and the induced (pullback) metric on the brane is
\begin{equation}
g_{ij}(x,t)=\partial_i\mathbf{X}\cdot\partial_j\mathbf{X}
=\sum_{A=1}^4 \partial_i X^A\,\partial_j X^A,\qquad i,j=1,\dots,d.
\label{eq:induced-metric}
\end{equation}

\paragraph{Reference (rest) metric and rest-length scale.}
A stress-free rest geometry is encoded by a reference metric $g^0_{ij}(x)$.
For a homogeneous isotropic rest state with a single rest-length scale $\ell_0>0$ we use
\begin{equation}
g^0_{ij}=\ell_0^2\,\delta_{ij},\qquad (g^0)^{ik}g^0_{kj}=\delta^{i}{}_{j}.
\label{eq:reference-metric}
\end{equation}
Define the mixed relative metric
\begin{equation}
\widehat g^{i}{}_{j}=(g^0)^{ik}g_{kj}.
\label{eq:relative-metric}
\end{equation}
Let $\lambda_a$ be the eigenvalues of $\widehat g^{i}{}_{j}$ and define relative principal stretches
\begin{equation}
\sigma_a=\sqrt{\lambda_a},\qquad a=1,\dots,d.
\label{eq:principal-stretches}
\end{equation}
The stress-free condition is $g=g^0$ (equivalently $\sigma_a=1$ for all $a$).

\paragraph{Stretch-only hyperelastic energy (no bending).}
We use an isotropic stretch-only energy density (per material volume) minimized at $g=g^0$:
\begin{equation}
W(g;g^0)=\frac{K}{2}\sum_{a=1}^{d}(\sigma_a-1)^2,
\qquad \sigma_a=\sqrt{\lambda_a(\widehat g)}.
\label{eq:stretch-energy-density}
\end{equation}
Here $K>0$ is the continuum stiffness scale and $W$ depends on $\mathbf{X}$ only through first derivatives
(via $g_{ij}$); no curvature/bending term is included in the present model statement.

\paragraph{Equations of motion.}
Let $\rho_m>0$ be mass density per material volume. The conservative equations of motion are
\begin{equation}
\rho_m\,\partial_{tt}X^A=\partial_i P_{Ai},\qquad A=1,\dots,4,
\label{eq:continuum-eom}
\end{equation}
with first Piola-type stress
\begin{equation}
P_{Ai}=\frac{\partial W}{\partial(\partial_i X^A)}
=2\,\frac{\partial W}{\partial g_{ij}}\,\partial_j X^A.
\label{eq:piola-stress}
\end{equation}
All four embedding components enter symmetrically through $\mathbf{X}$ and the induced metric $g_{ij}$; no
embedding axis is privileged by the dynamics.

\paragraph{Pre-tensioned base state.}
A pre-tensioned brane is defined by choosing a time-independent equilibrium embedding $\mathbf{X}_\star(x)$
such that its induced metric $g^\star_{ij}=\partial_i\mathbf{X}_\star\cdot\partial_j\mathbf{X}_\star$
differs from the reference metric: $g^\star\neq g^0$. A convenient uniform choice is an affine embedding
\begin{equation}
\mathbf{X}_\star(x)=A x+\mathbf b,\qquad A\in\mathbb{R}^{4\times d},\ \mathrm{rank}(A)=d,
\label{eq:affine-base-embedding}
\end{equation}
for which $g^\star=A^\mathsf{T}A$ is constant.
Pre-tension is maintained by boundary conditions or external loads (e.g.\ periodic cell with fixed shape, or
Dirichlet clamping $\mathbf{X}|_{\partial\Omega}=\mathbf{X}_\star|_{\partial\Omega}$). Perturbations are then
written as $\mathbf{X}(x,t)=\mathbf{X}_\star(x)+\boldsymbol\xi(x,t)$.

\paragraph{Dictionary (continuous model).}
\begin{itemize}
\item $d$: intrinsic/material dimension of the brane (typically $3$).
\item $\Omega\subset\mathbb{R}^d$: material domain; $x^i$ material coordinates; $t$ time parameter.
\item $\mathbf{X}(x,t)\in\mathbb{R}^4$: embedding map; $X^A$ its components ($A=1,\dots,4$).
\item $F_{Ai}=\partial_i X^A$: deformation gradient.
\item $g_{ij}=\partial_i\mathbf{X}\cdot\partial_j\mathbf{X}$: induced metric.
\item $g^0_{ij}=\ell_0^2\delta_{ij}$: reference (rest) metric; $\ell_0$ rest-length scale.
\item $\widehat g^{i}{}_{j}=(g^0)^{ik}g_{kj}$: relative metric; $\lambda_a$ its eigenvalues; $\sigma_a=\sqrt{\lambda_a}$ relative principal stretches.
\item $W(g;g^0)$: stretch-only energy density; $K$ stiffness scale.
\item $\rho_m$: mass density; $P_{Ai}$ first Piola-type stress.
\item $\mathbf{X}_\star$: pre-tensioned equilibrium embedding; $g^\star$ its induced metric.
\end{itemize}

\paragraph{Assumptions (continuous model).}
\begin{itemize}
\item Euclidean embedding space $\mathbb{R}^4$ with standard inner product.
\item No bending/curvature energy; $W$ depends only on first derivatives via $g_{ij}$.
\item Hyperelastic and isotropic material response: $W$ depends only on the eigenvalues of $(g^0)^{-1}g$.
\item Regular embedding in the regime of interest ($g_{ij}$ positive definite).
\item Pre-tension is a maintained equilibrium mismatch ($g^\star\neq g^0$) enforced by boundary conditions or loads.
\item No dissipation unless explicitly added in the implementation layer.
\end{itemize}

\subsubsection{Linear wave regime and isotropy}\label{linear-wave-regime-and-isotropy}

Linearize about a flat pre-tensioned embedding $\mathbf X_\star(x)=(x^1,x^2,x^3,0)$ with small displacement $\boldsymbol\xi(x,t)$ so $\mathbf X=\mathbf X_\star+\boldsymbol\xi$, and retain leading order. For a uniform pre-tension state with $g^\star_{ij}=\delta_{ij}$ and homogeneous parameters, one obtains a linear wave equation of the form
\begin{equation}
\partial_{tt}\boldsymbol\xi = c^2\Delta\boldsymbol\xi,
\label{eq:linear-wave-equation}
\end{equation}
where the effective wave speed is determined by the continuum stiffness and density:
\begin{equation}
c^2 = \frac{K}{\rho_m}.
\end{equation}
Throughout this work we identify this emergent wave speed with the physical speed of light $c$.

All \textbf{four components} of $\boldsymbol\xi$ obey the same \textbf{isotropic} wave equation in the material coordinates. The wave speed is the same for all directions and for all embedding coordinates. Geometric nonlinearities (large stretches) introduce deviations from this ideal isotropic behavior.

The relativistic interpretation of this linear wave regime, and its relation to
effective Lorentz symmetry, is developed in more detail in
Section~\ref{emergent-relativity-and-lorentz-invariance}.

\textbf{Regime of validity (boxed):}
The linear wave equation~\eqref{eq:linear-wave-equation} with isotropic
dispersion is valid in the small-strain regime, where geometric nonlinearities
can be neglected. In this regime the
wave speed $c=\sqrt{K/\rho_m}$ is approximately constant and independent of
direction. This linear, isotropic propagation behaviour underlies the emergent
Lorentz symmetry discussed in
Section~\ref{emergent-relativity-and-lorentz-invariance}.

\subsubsection{Geometric nonlinearity and localization threshold}\label{subsec:geom-threshold}

Define the local energy density
\begin{equation}
\mathcal{E}=\frac{\rho_m}{2}|\partial_t\mathbf{X}|^2
+W(g;g^0).
\label{eq:energy-density}
\end{equation}

\paragraph{Threshold for soliton nucleation: dimensionless nonlinearity.}
Geometric nonlinearity is controlled by the dimensionless parameter measuring the relative stretch:
\begin{equation}
\eta(x,t) \;\sim\; |\nabla \mathbf{X}|^2 \;\sim\; (A k)^2,
\label{eq:eta-nonlinearity}
\end{equation}
for a mode of amplitude $A$ and characteristic wavenumber $k$. Localization into a self-confined soliton occurs when $\eta$ reaches order unity over a region of linear size $\sim\lambda_C$, where $\lambda_C = \hbar/(m_e c)$ is the reduced Compton wavelength. For such a configuration with $k \sim 1/\lambda_C$, the nonlinearity condition $\eta \gtrsim 1$ implies $A \sim \lambda_C$.

Given the calibration $\rho_m \lambda_C^3 \approx m_e$ (see Section~\ref{amplitude-scale-calibration}), the elastic energy in a Compton-scale cell with $\eta \sim 1$ satisfies
\begin{equation}
E_{\text{cell}} \;\sim\; \rho_m\,A^2\,\omega_C^2\,\lambda_C^3
\;\sim\; \rho_m c^2 \lambda_C^3
\;\sim\; m_e c^2,
\label{eq:threshold-criterion}
\end{equation}
where $\omega_C = c/\lambda_C$ is the Compton angular frequency. Thus, the \emph{energy-in-a-Compton-volume} threshold is not a separate postulate but the energetic restatement of the geometric nonlinearity condition $\eta \gtrsim 1$ given the mass-density calibration. Above threshold, the geometric nonlinearity inherent in the full strain tensor can naturally lock in localized high-stretch 4D configurations, turning a propagating photon-like mode into a self-confined, particle-like soliton.

Earlier drafts included a toy intrinsic-curvature estimate for a localized bulge; it is omitted here to keep the focus on the localization threshold and the Berry/holonomy diagnostics.

\subsubsection{"Gravity" as effective metric}\label{gravity-as-effective-metric}

Internal signal propagation uses the instantaneous induced metric $g_{ij}=\partial_i\mathbf X\cdot\partial_j\mathbf X$. Local embedding deformations modify $g$, creating \textbf{geodesic focusing} for in-brane rays. We therefore identify \textbf{effective gravity} with curvature of $g$ induced by the embedding deformation. In practice we extract $g$ from simulation snapshots and ray-trace null-like paths in $(t,x)$ with interval
\begin{equation}
ds^2 = -c^2 dt^2 + g_{ij}(x,t)\,dx^i\, dx^j \quad\text{(with $x^i$ intrinsic/material coordinates)}.
\label{eq:effective-metric-interval}
\end{equation}

For numerical experiments we work on finite brane patches
$\Omega$ with physically motivated boundary conditions
(periodic, free, or clamped; see Sec.~\ref{experimental-setting}).

\subsection{From substrate waves to emergent gauge structure and holonomy}
\label{sec:brane-to-berry-to-em}

The central analytical step is this: \emph{a real-valued substrate can generate
a physically meaningful local $U(1)$ or $U(N)$ gauge structure}, but not by
interpreting raw FFT phases component-wise. Instead, the relevant gauge freedom
lives on the \emph{bundle of normalized complex mode envelopes after band
isolation} \cite{Berry1984,WilczekZee1984}.

\subsubsection{From real substrate to complex analytical signal}
\label{sec:spectral-layer}

The substrate displacement field $\boldsymbol\xi(x,t)=\mathbf X(x,t)-\mathbf X_0(x)$
is real-valued, $\boldsymbol\xi(x,t)\in\mathbb R^4$. To construct a complex
envelope with a well-defined local phase, we must first isolate a
\emph{positive-frequency component}. This can be done via:

\begin{enumerate}
\item \textbf{Analytic signal construction:} Apply a Hilbert transform (or
  equivalent filter) in time to retain only $\omega>0$ components, yielding a
  complex field $\xi_+(x,t)$ such that $\mathrm{Re}[\xi_+(x,t)]=\xi(x,t)$.

\item \textbf{Band-pass filter in $(\mathbf k,\omega)$ space:} Compute the
  space--time Fourier transform
  \begin{equation}
  \hat{\xi}^A(\mathbf k,\omega)
  =\int_{\mathbb R}\!dt\int_{\Omega}\!d^3x\;
  \xi^A(\mathbf x,t)\,e^{-i(\mathbf k\cdot\mathbf x-\omega t)},
  \qquad A=1,\dots,4,
  \label{eq:xi-fourier}
  \end{equation}
  and restrict to a narrow region in $(\mathbf k,\omega)$ space (a "band" in
  the broad sense: either a single eigen-branch $\omega_n(\mathbf k)$ or a
  localized wavepacket occupying a small tube in frequency space).
\end{enumerate}

In either case, the output is a complex four-component envelope field
$u(x,t)\in\mathbb C^4$ that varies slowly in space and time compared to its
internal oscillation scale.

\paragraph{Order of operations (transform-then-project).}
A central numerical subtlety is that \emph{geometric phase lives on a gauge bundle}:
the eigenvectors/polarization frames used for band identification can become
\emph{multi-valued} (sign flips / branch cuts) when the dynamics approaches
degeneracies or mode-crossings. In such situations, performing a Fourier/Wigner-like
transform \emph{after} an adiabatic projection can implicitly require transforming
double-valued objects and can suppress geometric-phase effects.
A clear example is given by the WA vs.\ AW comparison in quantum--classical Liouville
theory, where only the \emph{Wigner-then-Adiabatic} order captures GP effects
\cite{Ryabinkin2013QCLGP}. For the present brane programme we therefore adopt the
analogous rule:
\emph{when filtering is used as a diagnostic cross-check, apply any time--frequency or $(\mathbf{k},\omega)$ filtering to the
single-valued substrate fields first, and only then project onto the target band/subspace. Primary results, however, rely on
preparation-first narrowband initialization rather than post-hoc band-pass masks.}

\subsubsection{Band definition and gauge structure}
\label{sec:band-isolation-route2}

\paragraph{Preparation-first band definition (no diagnostic band filters).}
In this programme, the relevant ``band'' is defined by \emph{physical preparation}
of the substrate initial conditions: we initialize narrowband carrier wave packets
with a chosen $(\mathbf{k}_0,\omega_0)$ (typically anchored to the Compton scale)
and with a controlled internal polarization subspace. Diagnostic layers (FFT,
Berry, holonomy) are used only to \emph{read out} the resulting phase and
degeneracy structure. Spectral band-pass masks and dispersion-branch projections
are retained solely as secondary cross-checks and visualization aids.

\paragraph{Adiabatic representation (why eigenvalues and eigenvectors both matter).}
Linearizing the substrate dynamics around a slowly varying background yields a
parameter-dependent eigenproblem (schematically) $K(R)\,u_n(R)=\omega_n(R)^2\,M\,u_n(R)$,
where $R$ denotes the slowly varying parameters (e.g.\ $x$, $t$, or $\mathbf{k}$ after filtering).
In an \emph{adiabatic representation} the complex envelope is expanded in the instantaneous
eigenmodes, and the evolution separates into (i) a fast \emph{dynamical phase} controlled by the
eigenvalues $\omega_n$ and (ii) a \emph{geometric phase} controlled by the twisting of the eigenvectors
$u_n(R)$. The usual single-band ($U(1)$) Berry connection is well-defined only when the target branch
is isolated by a spectral gap; near degeneracies one must transport an $N$-dimensional subspace and
use the $U(N)$ (Wilczek--Zee) connection.

We use the term \emph{band} in a practical sense:

\paragraph{Definition (band).}
A \emph{band} is a smoothly varying family of complex mode directions in the
substrate's configuration space. In our primary workflow the band is defined by
\emph{preparation-first} narrowband initialization; the procedures below are used
mainly as diagnostic readout/cross-checks of that prepared state:
\begin{enumerate}
\item \textbf{Preparation-first initialization:} Initialize a narrowband carrier
  wave packet at a chosen $(\mathbf k_0,\omega_0)$ together with a controlled
  internal polarization subspace (possibly degenerate). This defines the target
  mode family \emph{before} any diagnostics are applied.

\item \textbf{Dispersion-branch isolation (cross-check):} For a locally translation-invariant
  substrate, compute the eigen-decomposition of the dynamical matrix
  $D(\mathbf k)$ and track a single branch $\omega_n(\mathbf k)$ with
  polarization eigenvector $u_n(\mathbf k)\in\mathbb C^4$.

\item \textbf{Wavepacket filtering (cross-check):} Select a region $\mathcal B\subset
  (\mathbf k,\omega)$ space (e.g., $|\omega-\omega_0|\leq\Delta\omega$,
  $|\mathbf k-\mathbf k_0|\leq\Delta k$) and inverse-transform only those
  Fourier components, yielding a band-limited envelope $u(x,t)$.

\item \textbf{Demodulation + low-pass filter (primary readout):} Multiply the substrate field by
  a carrier $e^{i(\mathbf k_0\cdot\mathbf x-\omega_0 t)}$ at the \emph{prepared}
  $(\mathbf k_0,\omega_0)$ and apply a spatial and temporal low-pass filter to
  extract the slowly varying envelope.
\end{enumerate}

In all three cases, the physically relevant object is the \emph{normalized}
complex direction
\begin{equation}
\tilde u(x,t) = \frac{u(x,t)}{|u(x,t)|},
\qquad
|u|^2 = \sum_{A=1}^{4} |u^A|^2.
\label{eq:normalized-mode}
\end{equation}

\paragraph{$U(1)$ gauge freedom.}
Although $u(x,t)\in\mathbb C^4$ has four complex components, the only gauge
redundancy is the overall phase:
\begin{equation}
\tilde u(x,t) \mapsto e^{i\chi(x,t)} \tilde u(x,t).
\end{equation}
The physically measurable content of the band is encoded in the projector
$P=|\tilde u\rangle\langle\tilde u|$, which is invariant under this $U(1)$ transformation.

\subsubsection{Connection and holonomy in spacetime}
\label{sec:berry-connection-curvature}

The primary definition of gauge structure in this framework is the
\emph{spacetime} connection one-form, defined directly from the normalized envelope:

\begin{equation}
a_\mu(x,t) = i \langle \tilde u(x,t), \partial_\mu \tilde u(x,t) \rangle
= i \sum_{A=1}^{4} \tilde u^{A\ast}(x,t) \, \partial_\mu \tilde u^A(x,t),
\label{eq:spacetime-berry-connection}
\end{equation}

\paragraph{Discrete, gauge-robust evaluation.}
While Eq.~\eqref{eq:spacetime-berry-connection} is the continuum definition,
numerical evaluation should avoid relying on raw derivatives of an arbitrarily
gauged eigenvector field. In practice we compute holonomy and curvature from
\emph{overlap products} (link variables / Wilson loops), which are manifestly
gauge-invariant and remain stable near gauge seams and (near-)degeneracies.
This choice is directly motivated by the fact that double-valued adiabatic objects
can make transform-based pipelines ill-defined and can erase GP physics if handled
in the wrong order \cite{Ryabinkin2013QCLGP}; see the implementation details in
Section~\ref{experimental-setting}.

where $\partial_\mu=(\tfrac{1}{c}\partial_t,\nabla)$ and $\langle\cdot,\cdot\rangle$
denotes the Hermitian inner product. This connection is real-valued and
gauge-invariant up to total derivatives: under $\tilde u\mapsto e^{i\chi}\tilde u$,
one has $a_\mu\mapsto a_\mu+\partial_\mu\chi$.

The associated curvature two-form is

\begin{equation}
f_{\mu\nu}(x,t) = \partial_\mu a_\nu - \partial_\nu a_\mu.
\label{eq:spacetime-berry-curvature}
\end{equation}

This curvature is fully gauge-invariant. The integral of $f_{\mu\nu}$ over a
closed two-surface gives the \emph{holonomy} (Berry phase) accumulated by
parallel-transporting the mode $\tilde u$ around the boundary of that surface.

\paragraph{Scalar-field caveat.}
If the band consists of a single scalar excitation (i.e., $u(x,t)$ is
effectively one-dimensional in the internal space), then the connection
$a_\mu=i\,\partial_\mu\log\tilde u$ and the curvature $f_{\mu\nu}$ vanishes
everywhere except at \emph{zeros} of $\tilde u$ (defect points/lines).
Smooth, non-zero curvature requires either:
\begin{enumerate}
\item a multi-component polarization structure (e.g., two orthogonal modes that
  vary in relative phase, giving nontrivial $\mathbb C^2$ geometry), or
\item degeneracies/defects where the band description breaks down.
\end{enumerate}

\subsubsection{Relation to k-space Berry connection (optional view)}
\label{sec:k-space-berry-optional}

For completeness, we note that in the case of dispersion-branch isolation
(method 1 above), one can alternatively compute the Berry connection in
$\mathbf k$-space:

\begin{equation}
\mathcal A_i(\mathbf k) = i\langle u_n(\mathbf k),\partial_{k_i}u_n(\mathbf k)\rangle,
\qquad i=1,2,3,
\label{eq:berry-connection-k}
\end{equation}

with curvature $\mathcal F_{ij}=\partial_{k_i}\mathcal A_j-\partial_{k_j}\mathcal A_i$.
If a wavepacket is localized around a slowly varying local wavevector
$\mathbf k(x,t)$, one can relate the spacetime and k-space connections via

\begin{equation}
a_\mu(x,t) \approx \mathcal A_i(\mathbf k(x,t))\,\partial_\mu k_i(x,t)
+ \text{(envelope-gradient corrections)}.
\label{eq:k-to-spacetime-pullback}
\end{equation}

However, this "pullback" picture is \emph{secondary}: it relies on the
semiclassical assumption that $\mathbf k(x,t)$ is well-defined, which fails
for strongly localized or defect-carrying configurations. The primary object is
the spacetime connection~\eqref{eq:spacetime-berry-connection}.

\subsubsection{Identification with effective gauge potential}
\label{sec:berry-pullback-em}

We define an \emph{effective} electromagnetic four-potential by introducing a
conversion factor (to be calibrated against the observed electron charge):

\begin{equation}
A_\mu^{\mathrm{eff}}(x,t) = \frac{\hbar}{q} \, a_\mu(x,t),
\label{eq:em-potential-from-berry}
\end{equation}

and the associated field-strength tensor

\begin{equation}
F_{\mu\nu}^{\mathrm{eff}}(x,t)
= \partial_\mu A_\nu^{\mathrm{eff}} - \partial_\nu A_\mu^{\mathrm{eff}}
= \frac{\hbar}{q} \, f_{\mu\nu}(x,t).
\label{eq:em-tensor-from-berry}
\end{equation}

By construction, $F_{\mu\nu}^{\mathrm{eff}}$ automatically satisfies the
\emph{homogeneous} Maxwell equations (Bianchi identities):

\begin{equation}
\partial_{[\lambda} F_{\mu\nu]}^{\mathrm{eff}} = 0,
\qquad \text{equivalently:} \quad
\nabla\cdot\mathbf B=0, \quad
\nabla\times\mathbf E+\partial_t\mathbf B=0.
\end{equation}

The \emph{inhomogeneous} Maxwell equations (dynamics with sources) are \textbf{not}
automatically satisfied; they become an empirical target:
\begin{itemize}
\item Do defects/solitons act as localized effective sources?
\item Does the long-wavelength envelope dynamics of band-projected modes
  reproduce wave equations consistent with Maxwell's equations in vacuum?
\end{itemize}

These questions are addressed numerically in Section~\ref{falsifiable-criteria}.

\paragraph{Interpretation: holonomy, not Maxwell.}
At this stage, we have defined a $U(1)$ connection and its curvature
geometrically. The physical content of this structure is:
\begin{enumerate}
\item \textbf{Holonomy / Aharonov--Bohm phase:} parallel transport of the mode
  $\tilde u$ around a closed loop accumulates a phase
  $\gamma=\oint a_\mu dx^\mu$, which can be measured via interference.

\item \textbf{Minimal coupling:} if an effective "matter" degree of freedom
  (e.g., the Dirac envelope of the electron soliton) couples to the band via
  $\partial_\mu\mapsto\partial_\mu-i(q/\hbar)A_\mu^{\mathrm{eff}}$, gauge
  invariance is preserved.
\end{enumerate}

We do \textbf{not} claim that this immediately \emph{is} electromagnetism in the
full Maxwell sense. That identification is a hypothesis to be tested via the
numerical experiments of Section~\ref{experimental-setting}.

\subsubsection{Degeneracies, defects, and non-Abelian holonomy}
\label{sec:nonabelian-berry-electron}

When the mode subspace is \emph{degenerate} (dimension $N>1$), the gauge
structure generalizes to the Wilczek--Zee non-Abelian connection
\cite{WilczekZee1984}:

\begin{equation}
\mathcal A_\mu^{ab}(x,t) = i \langle u_a(x,t), \partial_\mu u_b(x,t) \rangle,
\qquad a,b=1,\dots,N,
\end{equation}

where $\{u_1,\dots,u_N\}$ is an orthonormal frame spanning the degenerate
subspace, and the gauge group is $U(N)$.

In the present programme, robust localized defects (the toroidal electron
micro-structure described in Section~\ref{sec:internal-structure-of-the-electron})
are expected to act as degeneracy sources: the frame of polarization
eigenvectors becomes ill-defined at the defect core. Holonomy around such a
defect can yield a $U(2)\supset\mathrm{SU}(2)$ matrix, naturally encoding
spin-$1/2$ behavior (including $4\pi$ periodicity).

This supplies a concrete geometric mechanism for the electron's spinorial
character, distinct from the Abelian ($U(1)$) sector that describes the
electromagnetic interaction of charged matter.